% -----------------------------------------------------------------------------
% Chapter 11: Authentication & Authorization
% -----------------------------------------------------------------------------
\chapter{Authentication \& Authorization}
\label{chap:auth}

\section{Identity Propagation Pipeline}
Security within CreditTech is built on an end-to-end identity resolution pipeline. Unlike monolithic web applications that rely on opaque session cookies or client-side trust, CreditTech verifies caller identity at the gateway layer, derives an authenticated \texttt{CallerContext}, and passes this security context immutably through FastAPI dependency injection to downstream business logic. Figure~\ref{fig:identity_pipeline} illustrates this identity resolution flow.

\begin{figure}[H]
\centering
\begin{tikzpicture}[node distance=1.1cm and 1.2cm, auto]

  % Stage 1: Inbound Request
  \node[clientblock, fill=primary!10, draw=primary, minimum width=3.4cm] (req) {
    \textbf{Inbound HTTP Request}\\
    \scriptsize Headers: \texttt{X-Borrower-Id},\\
    \scriptsize \texttt{X-Officer-Role}, \texttt{X-Officer-Id}
  };

  % Stage 2: Gateway Interceptor
  \node[securityblock, right=1.0cm of req, fill=danger!10, draw=danger, minimum width=3.6cm] (gateway) {
    \textbf{Security Middleware}\\
    \scriptsize Rate Limit Check\\
    \scriptsize Body Size Sanitization\\
    \scriptsize Request-ID Ingestion
  };

  % Stage 3: Dependency Injection
  \node[serviceblock, right=1.0cm of gateway, fill=accent!10, draw=accent, minimum width=3.6cm] (dep) {
    \textbf{FastAPI Dependency}\\
    \scriptsize \texttt{require\_officer()}\\
    \scriptsize \texttt{require\_admin()}\\
    \scriptsize \texttt{get\_borrower\_ctx()}
  };

  % Stage 4: Resource Authorizer
  \node[decisionblock, below=1.3cm of dep, fill=warn!10, draw=warn, minimum width=3.2cm] (authz) {
    \textbf{Resource}\\
    \textbf{Authorizer}\\
    \scriptsize Tenant Match?\\
    \scriptsize Role Permitted?
  };

  % Stage 5: Execution vs Abort
  \node[baseblock, left=1.2cm of authz, fill=success!12, draw=success, minimum width=3.4cm] (exec) {
    \textbf{Execute Handler}\\
    \scriptsize Authorized Context\\
    \scriptsize Service Business Logic
  };

  \node[securityblock, below=1.1cm of authz, fill=danger!15, draw=danger, minimum width=3.4cm] (abort) {
    \textbf{Immediate Termination}\\
    \scriptsize HTTP 401 / 403 Abort\\
    \scriptsize Audit Alert Dispatched
  };

  % Connecting Flows
  \draw[flowline] (req) -- (gateway);
  \draw[flowline] (gateway) -- (dep);
  \draw[flowline] (dep) -- (authz);
  \draw[flowline] (authz) -- node[above, font=\scriptsize\sffamily] {Pass} (exec);
  \draw[flowline] (authz) -- node[right, font=\scriptsize\sffamily] {Fail} (abort);

\end{tikzpicture}
\caption{Research Block Diagram: Inbound Request Identity Propagation and Resource Authorization Pipeline.}
\label{fig:identity_pipeline}
\end{figure}

\section{Defense Against Horizontal Privilege Escalation}
Horizontal privilege escalation occurs when an attacker possesses valid credentials as a regular user (e.g., Borrower A) and attempts to access or manipulate private resources belonging to another user of equivalent rank (e.g., Borrower B).

In CreditTech, this threat is neutralized at the route dependency layer through an explicit \textit{Tenant Ownership Invariant}:
\begin{lstlisting}[language=Python, caption={Tenant Ownership Verification in Score Retrieval Handler.}]
@router.get("/score/{score_id}", summary="Fetch credit score and SHAP reason codes")
async def get_score(
    score_id: UUID,
    borrower_id: UUID | None = Depends(get_optional_borrower_id),
    officer: OfficerContext | None = Depends(get_optional_officer),
    db: AsyncSession = Depends(get_db),
):
    score = await db.get(ScoreResult, score_id)
    if not score:
        raise HTTPException(status_code=404, detail="Score record not found")
        
    # Horizontal Authorization Guard:
    if borrower_id is not None:
        # Caller is a borrower -> Must own the associated snapshot
        snapshot = await db.get(FeatureSnapshot, score.feature_id)
        if snapshot.borrower_id != borrower_id:
            logger.warn("horizontal_privilege_attempt", caller=borrower_id, target=snapshot.borrower_id)
            raise HTTPException(status_code=403, detail="Access to resource forbidden: caller does not own record")
            
    return score.to_response()
\end{lstlisting}

\section{Defense Against Vertical Privilege Escalation}
Vertical privilege escalation occurs when a standard user (e.g., a Borrower or Field Sakhi) attempts to execute restricted operations reserved for higher-privilege administrative or underwriting roles (e.g., activating an unvetted ML model or overriding an adverse loan decision).

CreditTech prevents vertical escalation through tiered role-guard dependencies (\texttt{require\_officer}, \texttt{require\_admin}). Any route modifying platform-level configuration or model promotion enforces strict role membership:
\begin{lstlisting}[language=Python, caption={Role-Gating Dependency Implementation.}]
async def require_admin(
    role: str = Header(..., alias="X-Officer-Role"),
    officer_id: str = Header(..., alias="X-Officer-Id"),
) -> str:
    if role.strip().upper() != "ADMIN":
        logger.error("vertical_escalation_blocked", caller=officer_id, attempted_role=role)
        raise HTTPException(
            status_code=status.HTTP_403_FORBIDDEN,
            detail="Operation requires administrative privileges"
        )
    return officer_id
\end{lstlisting}
Because these checks are evaluated before handler execution, no sensitive internal state, weights, or database queries are leaked to unauthorized callers.
